The interpreter's cost is its run time on the CPU. \figref{fig:overhead} and \tabref{tab:overhead} compare it with
the program's own run on the GPU and with the run under the NVBit tool, for six programs at three input sizes
each. The GPU runs take \numNativeLo{} to \numNativeHi{} seconds, and almost all of that time is the start-up of
the CUDA runtime, because at these sizes the kernels themselves take milliseconds. Against the whole GPU run, the
interpreter in lockstep takes \numOverheadLo{} to \numOverheadHi{} times as long, and the NVBit tool \numNvbitLo{}
to \numNvbitHi{} times. Against the kernels alone, the interpreter is several orders of magnitude slower, as
expected from simulating every thread in Python on one CPU core.

\textit{Most of the interpreter's time is spent in host code, not in kernels.} The ``In kernels'' column of
\tabref{tab:overhead} gives the share of the time spent running kernels. The kernels of \code{chain}, with
\numChainThreads{} GPU threads, take \numChainKernelS{} seconds, and those of \code{rodinia\_bfs}, with
\numRbfsThreads{} threads, take less than one second, so \code{rodinia\_bfs} as a whole takes \numRbfsLockS{}
seconds, against \numRbfsNvbitS{} seconds under NVBit. The time spent in kernels alone is \numKernelRatioLo{} to
\numKernelRatioHi{} times the whole GPU run, with a median of \numKernelRatioMedian{}. The programs that take
longest, \code{wl\_bfs} (\numWlbfsLockS{} seconds, of which \numWlbfsKernelS{} in kernels) and \code{bfs}
(\numBfsLockS{} seconds, of which \numBfsKernelS{} in kernels), spend up to \numHostShareMax{} of their time in host
loops that carry a value from one iteration to the next. These loops are the random-number generator that builds
the graph in \code{wl\_bfs} and \code{cc}, and a running maximum over the result in \code{bfs}. The interpreter runs
all iterations of a loop at once only when the iterations do not depend on each other, so it runs these loops one
iteration at a time, at about a millisecond per iteration. Such loops belong to the test harness here, not to the
computation on the GPU, but real programs contain them as well, for example when they read or generate their input.

Running all four thread orders costs \numScheduleLo{} to \numScheduleHi{} times as much as lockstep alone
(\numScheduleCost{} on average). Where host code dominates, the cost is close to four times, because each order runs
the host code once more. Where kernels dominate, the orders that run one block at a time cost much more than
lockstep, because they carry out every statement once per block instead of once per group of up to $2^{18}$ threads.
A tool could run the extra orders only for kernels that read memory that other threads of the same launch write,
which is where the order can change the result. We have not done so.
